\documentclass[11pt,a4paper]{article}

% ─── Packages ──────────────────────────────────────────────────────────────────
\usepackage[utf8]{inputenc}
\usepackage[T1]{fontenc}
\usepackage{amsmath,amssymb}
\usepackage{graphicx}
\usepackage{booktabs}
\usepackage{hyperref}
\usepackage{natbib}
\usepackage[margin=1in]{geometry}
\usepackage{algorithm}
\usepackage{algpseudocode}
\usepackage{xcolor}
\usepackage{enumitem}
\usepackage{microtype}
\usepackage{pifont}
\usepackage{multirow}

\hypersetup{
    colorlinks=true,
    linkcolor=blue,
    citecolor=blue,
    urlcolor=blue,
    hypertexnames=false,
}

% ─── Title ─────────────────────────────────────────────────────────────────────

\title{The Retrieval--Reasoning Gap: Why Better Retrieval Does Not Guarantee\\Better Reasoning in Iterative RAG Systems}

\author{Anugna Yakkala \\[4pt]
{\normalsize School of Computer Science and Engineering,}\\
{\normalsize Vellore Institute of Technology, Vellore, TN, India}}

\date{}

\begin{document}
\maketitle

% ═══════════════════════════════════════════════════════════════════════════════
% ABSTRACT
% ═══════════════════════════════════════════════════════════════════════════════

\begin{abstract}
Typical evaluations conducted using RAG evaluate the retrieval
itself and infer that any improvements in retrieval will provide
an improvement in the performance of the end-use task; however,
we have demonstrated this to be incorrect for multi-hop reasoning
tasks.
Specifically, we have established the \emph{retrieval--reasoning gap}
that exists between these two components: improvements made in
metrics on the retrieval side do not correlate with improved
accuracy of reasoning performance when the language model is fixed.
To investigate this further, an ablation study with six
configurations on a set of 65 multi-hop queries created from
synthetic telecom logs spanning a total of 1{,}115 records has
been performed.
Within this analysis, we found that the Precision@K values for
this ablation set would vary by an average of 0.24 ($p{<}0.001$),
but that the root-cause accuracy for the same dataset had no
statistically significant mean change ($\Delta{=}0.05$, $p{=}0.08$,
n.s.; Pearson $r{=}0.244$).
Additionally, we conducted experiments using multiple variants of
the same LLM (8--70 billion parameters) and confirmed that while
retrieval is impacted by model capability, the accuracy of
reasoning remains constant across configurations once retrieval
has been set to a particular value.
To assess the diagnostics of our models against one another, we
constructed a composite metric for evaluating performance based on
semantic similarity, ROUGE-L, and structured scoring to capture
the diagnostic capabilities of our model going beyond standard
word overlap ratios ($r{=}0.849$ with semantic evaluation).
Our data do generalize to multimodal pipelines; any system that
uses heterogeneously assembled data as inputs to a generative
process after retrieving relevant evidence will have the same
structural gap issues related to retrieval and reasoning.
\end{abstract}

% ═══════════════════════════════════════════════════════════════════════════════
% 1. INTRODUCTION
% ═══════════════════════════════════════════════════════════════════════════════

\section{Introduction}

The Retrieval-Augmented Generation (RAG) systems
\citep{lewis2020rag,gao2023ragsurvey} are evaluated according to
two families of metrics: retrieval quality (Precision@K, Recall@K,
nDCG) and end-task quality (accuracy of answer, F1).
Implicitly, we assume that improvements to retrieval lead to
improvements to the end-task.
The assumption of correlation for multi-hop retrieval-augmented
generation is an important but under-studied assumption since
multi-hop RAG requires the assembly of evidence from multiple
sources prior to the reasoning of multiple associations by an LLM.

In this paper, we challenge our core assumption and back it with
controlled empirical evidence that shows that improving retrieval
on its own can be insufficient as the reader model for multi-hop
is constrained at all times.

We present the \textbf{retrieval--reasoning gap}: an observation
that improvements to retrieval-side metrics do not translate to
improvements in the accuracy of reasoning, and that the
retrieval--reasoning gap is a structural property of multi-hop
retrieval-augmented generation and not a function of the size or
type of dataset being evaluated.

To isolate retrieval from reasoning in our experiments, we use
three approaches:
\begin{enumerate}[nosep]
    \item An \textbf{ablation study using 6 configurations} of the
    retrieval pipeline ($\pm$\,BM25, $\pm$\,cross-encoder reranking,
    $\pm$\,iterative refinement) with the LLM model held constant
    (Section~\ref{sec:ablation}).
    \item A \textbf{multi-LLM comparison} with fixed retrieval and
    only changing the reasoning model across four unique architectures
    of 8B--70B parameters (Section~\ref{sec:multi-llm}).
    \item \textbf{New evaluation metrics}---semantic similarity,
    structured scoring, composite metric---that assess reasoning
    quality beyond the degree of keyword overlap
    (Section~\ref{sec:metrics-result}).
\end{enumerate}

We performed our evaluation on a synthetic telecom log dataset
consisting of 1{,}115 records, 10 log files and 10 distinct failure
scenarios, and 65 evaluation queries requiring multi-hop reasoning
across files.
While the telecom domain serves as a controlled testbed, the
evaluation phenomena we discuss---including metric disagreements,
independence of failure modes, and inadequacy of aggregate
scores---all apply at a general level for any RAG system performing
multi-hop evidence assembly followed by LLM-based reasoning.

Our experiments were conducted in a text-based modality context, but
the retrieval--reasoning gap is not modality-dependent.
Multimodal retrieval systems that retrieve heterogeneous evidence
across multiple modalities (image+text, video+text, structured data)
face the same fundamental challenge in that retrieval of relevant
support materials does not necessitate downstream multimodal
generator(s) reasoning correctly with that association.
Thus, our proposed evaluation framework of decomposing retrieval
from reasoning applies directly to these types of multimodal systems.

\paragraph{Contributions.}
\begin{enumerate}[nosep]
    \item \textbf{Identification that the retrieval--reasoning gap is
    a structural property.}
    Across the six ablation configurations, Precision@K showed a 0.24
    variation ($p{<}0.001$) while the accuracy associated with the
    root cause showed only a 0.05 change ($p{=}0.08$, n.s.;
    Pearson $r{=}0.244$).
    In addition, the multi-LLM comparison with fixed retrieval
    demonstrates up to a 0.103 gap in reasoning accuracy attributable
    solely to the capability of the model.
    \item \textbf{Decomposition protocol for multi-hop RAG evaluation.}
    We propose: (a)~ablation with fixed LLM, (b)~model comparison
    with fixed retrieval, and (c)~per-query regime classification of
    high/low retrieval $\times$ correct/incorrect reasoning---a
    framework applicable to both text-only and multimodal approaches.
    \item \textbf{Composite evaluation metric.}
    We show that the degree of keyword overlap is a noisy measure of
    reasoning quality and propose a composite score (semantic
    similarity + ROUGE-L + structured scoring) for capturing
    diagnostic quality to a greater degree.
    \item \textbf{Dataset release.}
    Synthetic telecom log dataset: 1{,}115 records across 10 failure
    scenarios and 65 annotated queries with associated ground-truth
    root causes.
\end{enumerate}

% ═══════════════════════════════════════════════════════════════════════════════
% 2. RELATED WORK
% ═══════════════════════════════════════════════════════════════════════════════

\section{Related Work}
\label{sec:related}

\paragraph{Retrieval-Augmented Generation.}
Retrieval-Augmented Generation (RAG), with \citet{lewis2020rag}
providing empirical evidence that extending LMs with retrieved
passages provides large benefits with knowledge-intensive tasks, and
\citet{gao2023ragsurvey} summarizing extensions to various
summarization, code generation, and domain-specific applications.
However, the majority of RAGs utilize a single-pass architecture,
inadequate for analyzing multi-file log files.

\paragraph{Iterative and Multi-Step Retrieval.}
The use of self-reflective retrieval with Self-RAG
\citep{asai2024selfrag}, the ability to utilize active learning
signals in Active RAG \citep{jiang2023activerar}, and the use of
query rewriting to reformulate queries after multiple passes
\citep{wang2023queryrewrite} have not provided stop criteria.
Our work addresses the lack of principled stop criteria through our
adaptive confidence-gated iteration, and our ablation study
quantifies the benefit of iteration for both retrieval and reasoning.

\paragraph{Hybrid Retrieval and Cross-Encoder Reranking.}
Using BM25 \citep{robertson2009bm25} and a bi-encoder retrieval
system \citep{reimers2019sentencebert} is considered a standard
practice in the RAG community.
Studies show that using a cross-encoder reranker
\citep{nogueira2019passage} significantly improves the precision of
retrieved records and BEIR \citep{thakur2021beir} has confirmed the
ability of a cross-encoder reranker to provide zero-shot
generalization.
Our analysis shows that while reranking does provide substantial
benefits to precision of retrieval, it does not proportionally affect
the accuracy of reasoning.

\paragraph{RAG Evaluation.}
The RAGAS \citep{es2024ragas} and ARES \citep{saad2024ares}
evaluation frameworks decompose the quality of RAGs into three parts:
context relevance, faithfulness to the source, and relevance of the
answer.
Our work contributes to this literature by illustrating that
retrieval and end-task metrics move independently of one another, and
neither metric alone or together can provide evidence as to which of
the reported failure modes was responsible for the decline in
performance.

% ═══════════════════════════════════════════════════════════════════════════════
% 3. METHOD
% ═══════════════════════════════════════════════════════════════════════════════

\section{Method}
\label{sec:method}

\subsection{Pipeline Architecture}

The four-stage iterative feedback loop consists of:
\emph{Hybrid Candidate Generation $\to$ Cross-Encoder Reranking
$\to$ LLM Analysis $\to$ Confidence-Gated Iteration}.

\paragraph{Stage 1: Hybrid Candidate Generation.}
The first stage is a hybrid candidate generation system that uses a
dense retriever (\texttt{all-MiniLM-L6-v2} indexed in ChromaDB using
384-dimensional embeddings) and BM25 (with parameters $k_1{=}1.5$
and $b{=}0.75$) in a weighted manner:
$s_{\text{hybrid}}(d) = 0.7 \cdot s_{\text{dense}}(d) + 0.3 \cdot
s_{\text{BM25}}(d)$,
resulting in 20 to 30 candidates for the second stage of
cross-encoder reranking.

\paragraph{Stage 2: Cross-Encoder Reranking.}
The \texttt{ms-marco-MiniLM-L-6-v2} model uses cross-attention to
score each (query, document) pair, and the assigned score is the
only way to rank the candidates, so any hybrid scoring (used in the
first stage) is ignored to mitigate any return of noise from the
first stage.
The top 6 scores of the candidates move to the next stage for
analysis.

\paragraph{Stage 3: LLM Analysis.}
The LLM (default: Llama~3.3~70B model, temperature~0) creates the
structured output of root cause, severity of issue, resolution time
frame, logical steps in the reasoning process, and an actionable
recommendation.
The final confidence that identifies which RAG iteration is to be
reported is derived from a composite of: 30\% of the mean score of
the retrieval stage, 30\% based upon the completeness of the output,
20\% based upon the density of supporting evidence of the output,
20\% based on the RAG iterations having similar output values.

\paragraph{Stage 4: Confidence-Gated Iteration.}
The confidence-gated RAG system has three distinct stopping rules
that will end an iteration cycle: exit if the combined overall
confidence is equal to or greater than 0.85, if either of the last
two RAG iterations had no improvement (a plateau), or if the
iteration cycle has exceeded three (maximum iterations permitted) to
identify the best-scoring iteration for the RAG system.

\begin{algorithm}[t]
\caption{Adaptive Iterative RAG}
\label{alg:adaptive}
\begin{algorithmic}[1]
\Require Query $q$, documents $\mathcal{D}$, max iters $M$, threshold $\epsilon$
\State $q_1 \gets q$; $c_0 \gets 0$; $\text{best} \gets \textsc{null}$; $\text{seen} \gets \emptyset$
\For{$i = 1$ \textbf{to} $M$}
    \State $\mathcal{C}_i \gets \textsc{HybridRetrieve}(q_i, \mathcal{D}, N_i)$
    \State $\mathcal{R}_i \gets \textsc{Rerank}(q_i, \mathcal{C}_i, k)$
    \State $\mathcal{R}_i \gets \textsc{Dedup}(\mathcal{R}_i, \text{seen})$; update $\text{seen}$
    \State $a_i \gets \textsc{LLM}(q_i, \mathcal{R}_i)$
    \State $c_i \gets \textsc{Confidence}(a_i, \mathcal{R}_i, a_{i-1})$
    \If{$c_i > c_{\text{best}}$} $\text{best} \gets (a_i, \mathcal{R}_i, c_i)$ \EndIf
    \If{$c_i \geq 0.85$ \textbf{or} drops $\geq 2$} \textbf{break} \EndIf
    \State $q_{i+1} \gets \textsc{RefineQuery}(q, q_i, a_i)$
\EndFor
\State \Return $\text{best}$
\end{algorithmic}
\end{algorithm}

\subsection{Implementation Details}

Retrieval Model: \texttt{all-MiniLM-L6-v2} (22M parameters,
384-dim embeddings) indexed in ChromaDB (in-memory).
Cross-Encoder Reranking Model: \texttt{ms-marco-MiniLM-L-6-v2}
(22M parameters, trained against the MS~MARCO passage dataset).
LLM (default): Llama~3.3~70B via Groq (temperature set to~0,
\texttt{max\_tokens=2048}).
Retrieval and cross-encoder reranking occurs during each session and
uses a CPU for computation to generate embeddings and perform
subsequent reranking.
The LLM model is called via an API.
All model types used in the adaptive iterative RAG system are
implemented without domain-specific fine-tuning or adjustment.

% ═══════════════════════════════════════════════════════════════════════════════
% 4. EXPERIMENTAL SETUP
% ═══════════════════════════════════════════════════════════════════════════════

\section{Experimental Setup}
\label{sec:setup}

\subsection{Dataset}

We created a synthetic telecom log collection to evaluate multi-hop
retrieval: there are \textbf{1{,}115 logs} across \textbf{10 log
files} from 6 different network domains (e.g., eNodeB, CU-CP, CU-UP,
RAN, Core Network, Transport).
The logs represent \textbf{10 different failure scenarios} and consist
of a chain of cause-and-effect relationships across 2--5 log files:

\begin{enumerate}[nosep]
    \item RRC Reconfiguration Failure --- 4 log files, 13 steps
    \item SCTP Transport Link Failure --- 45 UEs impacted (4 log files)
    \item Inter-gNB Handover Failure with PDCP Data Loss (4 log files)
    \item PDU Session Denied by PCF Policy (5 log files)
    \item Max RLC Retransmission Causes DRB Failure (4 log files)
    \item GTP-U Path Failure on Backhaul (4 log files)
    \item 5G-AKA Authentication Failure with SQN Desynchronization (4 log files)
    \item Downlink Scheduler Starvation Under Load (3 log files)
    \item IPSec Tunnel Failure on Midhaul (3 log files)
    \item AMF Overloaded by Registration Storm (3 log files)
\end{enumerate}

Each log file includes a series of logs that outline a scenario and
80--120 logs that document normal operations.
The user must retrieve the logs that provide relevant evidence for
their query.

\paragraph{Dataset Rationale.}
The logs for our synthetic dataset do not contain personally
identifiable information (PII) or private proprietary information
found in commercial-grade telecom logs (for example,
subscriber-specific configuration information).
A synthetic corpus has three benefits for this research:
(i)~The researcher has complete control over the complexity of the
failures presented in a corpus (i.e., type of failure and number of
causal relationships), allowing consistent evaluation of the efficacy
of the various methods used to identify root causes;
(ii)~every query in this study has a known root cause, allowing
researchers to avoid confusion when annotating the logs; and
(iii)~all of the data and code will be released upon acceptance
for replication of any results that are reported.

\subsection{Queries}

\textbf{65 evaluation queries} were included as follows:
root-cause analysis (10), failure tracing (10), multi-hop (10),
temporal (10), impact analysis (10), error code (10),
cross-scenario (4), contrastive (1).
Of the total number of queries submitted, 55 required cross-file
reasoning to determine the correct retrieval answer.
Each query was annotated for the log file(s), root-cause keyword(s),
and free text describing the true nature of the logs related to
the query.

\subsection{Ablation Configurations}

There were six configurations of the models where each model used
LLM (Llama~3.3~70B):

\begin{table}[h]
\centering
\small
\caption{Ablation configurations.}
\label{tab:configs}
\begin{tabular}{lcccc}
\toprule
\textbf{Config} & \textbf{Dense} & \textbf{BM25} & \textbf{Rerank} & \textbf{Iter.} \\
\midrule
Dense-Only       & \checkmark &            &             &             \\
BM25-Only        &            & \checkmark &             &             \\
Hybrid           & \checkmark & \checkmark &             &             \\
Hybrid+Rerank    & \checkmark & \checkmark & \checkmark  &             \\
Hybrid+Iter.     & \checkmark & \checkmark &             & \checkmark  \\
Full System      & \checkmark & \checkmark & \checkmark  & \checkmark  \\
\bottomrule
\end{tabular}
\end{table}

\subsection{Multi-LLM Comparison}

All retrievals were fixed (Hybrid+Rerank, top-6) across four models:
Llama~3.3~70B, Llama~3.1~8B, Mixtral~8x7B, Gemma2~9B (all used
Groq; temperature~$=$~0).

\subsection{Metrics}

\textbf{Retrieval}: Precision@K, Recall@K, Mean Reciprocal Rank (MRR).
\textbf{Reasoning}: keyword overlap (RCA\textsubscript{kw}),
semantic similarity, structured scores, composite score
($0.35{\cdot}\text{sem} + 0.30{\cdot}\text{struct} +
0.20{\cdot}\text{ROUGE-L} + 0.15{\cdot}\text{kw}$).
\textbf{LLM-as-judge}: correct/not correct (binary).

% ═══════════════════════════════════════════════════════════════════════════════
% 5. RESULTS
% ═══════════════════════════════════════════════════════════════════════════════

\section{Results}
\label{sec:results}

\subsection{Study of Individual Components}
\label{sec:ablation}

There were multiple important findings for the aggregate measures
(Table~\ref{tab:ablation}).

\begin{table}[h]
\centering
\small
\caption{Ablation results (65 queries, Llama 3.3 70B).
Values produced by \texttt{run\_ablation.py} and filled at run time.}
\label{tab:ablation}
\begin{tabular}{lcccccc}
\toprule
\textbf{Config} & \textbf{P@K} & \textbf{R@K} & \textbf{MRR} & \textbf{RCA\textsubscript{kw}} & \textbf{Conf.} & \textbf{Iters} \\
\midrule
Dense-Only           & 0.553$\pm$0.18 & 0.500$\pm$0.19 & 0.613$\pm$0.18 & 0.311$\pm$0.16 & 0.553 & 1.0 \\
BM25-Only            & 0.598$\pm$0.18 & 0.509$\pm$0.21 & 0.629$\pm$0.17 & 0.314$\pm$0.16 & 0.550 & 1.0 \\
Hybrid               & 0.699$\pm$0.16 & 0.637$\pm$0.19 & 0.757$\pm$0.14 & 0.312$\pm$0.17 & 0.588 & 1.0 \\
Hybrid+Rerank        & 0.794$\pm$0.15 & 0.720$\pm$0.17 & 0.822$\pm$0.14 & 0.344$\pm$0.16 & 0.641 & 1.0 \\
Hybrid+Iter.         & 0.628$\pm$0.19 & 0.704$\pm$0.16 & 0.716$\pm$0.16 & 0.338$\pm$0.15 & 0.606 & 2.0 \\
Full-System          & 0.751$\pm$0.15 & 0.729$\pm$0.15 & 0.811$\pm$0.14 & 0.359$\pm$0.18 & 0.649 & 2.2 \\
\bottomrule
\end{tabular}
\end{table}

\paragraph{Findings.}
The most significant finding was:
(i)~A cross-encoder reranking produced the largest single-component
improvement: Hybrid$\to$Hybrid+Rerank produced a change in P@K of
+0.096 ($p{<}0.001$); however, there was only a small additional
increase of +0.032 ($p{=}0.234$, n.s.) for
RCA\textsubscript{kw}.
(ii)~A BM25 fusion produced additional useful information to support
the use of an exact error code when compared to the performance from
using the Dense-Only model alone and increased the P@K from 0.553
(Dense-Only) to 0.699 (Hybrid).
(iii)~Iteration without any reranking will produce lower precision
scores because of query drift (Hybrid+Iter.\
P@K$=$0.628 vs.\ Hybrid 0.699); conversely, it produced a higher
recall score (0.704 vs.\ 0.637).
(iv)~In all six configurations tested, the average
retrieval--reasoning gap for \textbf{P@K was 0.24 ($p{<}0.001$);
however, the average retrieval--reasoning gap for
RCA\textsubscript{kw} was only 0.05 ($p{=}0.08$, n.s.)}.
Therefore, this provides direct empirical evidence for the
retrieval--reasoning gap.
There are scatter plots of retrieval vs.\ reasoning accuracy
ratings available in the supplementary data section that include all
390 combinations of query$\times$configuration.

\subsection{Comparison of Multiple LLMs}
\label{sec:multi-llm}

\begin{table}[h]
\centering
\small
\caption{Multi-LLM results with identical retrieval.
Values from \texttt{run\_multi\_llm\_comparison.py}.}
\label{tab:multi-llm}
\begin{tabular}{lcccc}
\toprule
\textbf{Model} & \textbf{RCA\textsubscript{kw}} & \textbf{Completeness} & \textbf{Latency} & \textbf{Params} \\
\midrule
Llama-3.3-70B        & 0.363$\pm$0.14 & 0.703$\pm$0.10 & 4.4 & 70B \\
Mixtral-8x7B         & 0.300$\pm$0.12 & 0.674$\pm$0.09 & 3.8 & 46.7B \\
Gemma2-9B            & 0.302$\pm$0.13 & 0.632$\pm$0.10 & 2.1 & 9B \\
Llama-3.1-8B         & 0.260$\pm$0.12 & 0.607$\pm$0.12 & 1.9 & 8B \\
\bottomrule
\end{tabular}
\end{table}

The average P@K rating was nearly the same (P@K$\approx$0.794 for
all LLMs) when retrieval using the Hybrid+Rerank model was applied
to each of the four proposed models; however, the corresponding
RCA\textsubscript{kw} ranged from 0.260 (Llama~3.1~8B) to 0.363
(Llama~3.3~70B): a range of 0.103, which was due strictly to the
model capability difference.
Completeness followed the same pattern as for the variance by LLM
(0.607--0.703), thereby confirming \textbf{a model capability
independency basis as an influencing factor in reasoning success}.
Additionally, the overall drop from the 70B$\to$8B LLM in terms of
RCA due to reasoning dropped more (0.103) than the
Dense-Only$\to$Hybrid+Rerank retrieval-driven RCA change (0.032),
indicating that the total effect of retrieval optimization on
reasoning quality was dominated by LLM scaling.

\subsection{Correlation Between Retrieval and Reasoning}

\begin{table}[h]
\centering
\small
\caption{Correlation between retrieval and reasoning across all runs.
Values from \texttt{run\_correlation\_analysis.py}.}
\label{tab:correlation}
\begin{tabular}{lcccc}
\toprule
\textbf{Metric} & \textbf{Pearson $r$} & \textbf{$p$} & \textbf{Spearman $\rho$} & \textbf{$p$} \\
\midrule
Precision@K  & 0.244 & $<$0.001 & 0.222 & $<$0.001 \\
Recall@K     & 0.208 & $<$0.001 & 0.178 & 0.0004 \\
MRR          & 0.169 & 0.0008 & 0.143 & 0.0045 \\
\bottomrule
\end{tabular}
\end{table}

All correlations between retrieval metrics and reasoning accuracy
metrics are weak ($|r|{<}0.25$, $|\rho|{<}0.23$) but statistically
significant ($p{<}0.005$, which indicates that this is due to the
size of the sample).
The Pearson $r{=}0.244$ for P@K resulting in a retrieval quality
variance of approximately 6\% (out of a maximum 100\% variance)
being explained by the retrieval quality; thus \textbf{confirming the
existence of a retrieval--reasoning gap for our empirical data}.

\subsection{Query Failure Distributions}

\begin{table}[h]
\centering
\small
\caption{Per-query regime classification.}
\label{tab:failure-modes}
\begin{tabular}{lcc}
\toprule
\textbf{Regime} & \textbf{Count} & \textbf{\%} \\
\midrule
High retrieval, high reasoning           & 111 & 28.5 \\
High retrieval, low reasoning            & 84 & 21.5 \\
Low retrieval, high reasoning            & 84 & 21.5 \\
Low retrieval, low reasoning             & 111 & 28.5 \\
\bottomrule
\end{tabular}
\end{table}

43\% of queries (21.5\% High~P@K/Low~RCA + 21.5\%
Low~P@K/High~RCA) fall into off-diagonal quadrants---direct
evidence that aggregate-level metrics do not allow for
characterization of multiple failure modes.
Also, the symmetric distribution of the overall failure patterns
between the quadrants supports the independence claim since the
success of retrieval does not ensure or deny correct reasoning.

\subsection{Statistical Comparison of Findings}

Table~\ref{tab:significance} provides results of paired $t$-tests
on matched sample ($n{=}65$) queries with bootstrapped 95\%
confidence intervals.

\begin{table}[h]
\centering
\small
\caption{Statistical significance of key comparisons (paired $t$-test, $n{=}65$).}
\label{tab:significance}
\begin{tabular}{llccc}
\toprule
\textbf{Comparison} & \textbf{Metric} & \textbf{$\Delta$} & \textbf{$p$} & \textbf{95\% CI} \\
\midrule
Hybrid $\to$ Hybrid+Rerank & P@K & +0.096 & $<$0.001*** & [0.057, 0.134] \\
Hybrid $\to$ Hybrid+Rerank & RCA & +0.032 & 0.234 & [$-$0.021, 0.085] \\
\addlinespace
Dense-Only $\to$ Full System & P@K & +0.198 & $<$0.001*** & [0.149, 0.247] \\
Dense-Only $\to$ Full System & RCA & +0.048 & 0.083 & [$-$0.006, 0.102] \\
\addlinespace
Hybrid $\to$ Full System & P@K & +0.053 & 0.009** & [0.013, 0.092] \\
Hybrid $\to$ Full System & RCA & +0.047 & 0.132 & [$-$0.014, 0.107] \\
\bottomrule
\end{tabular}
\end{table}

Retrieval gained statistically significant improvements
($p{<}0.01$); reasoning differences were consistently
non-significant ($p{>}0.08$).
This observed pattern is representative of the statistical signature
that demonstrates the presence of a retrieval--reasoning gap for our
empirical data.

\subsection{Enhanced Metric Analysis}
\label{sec:metrics-result}

\begin{table}[h]
\centering
\small
\caption{Metric inter-correlation.
Values from \texttt{run\_improved\_metrics.py}.}
\label{tab:metric-corr}
\begin{tabular}{lcc}
\toprule
\textbf{Metric Pair} & \textbf{Pearson $r$} \\
\midrule
Keyword overlap vs.\ Semantic similarity      & 0.673 \\
Keyword overlap vs.\ Composite score          & 0.773 \\
Semantic sim.\ vs.\ ROUGE-L F1                & 0.628 \\
\bottomrule
\end{tabular}
\end{table}

The association between keyword overlap and semantic similarity is
only moderate ($r{=}0.673$), thereby indicating this metric is most
appropriate for use as a noisy surrogate measure; therefore there
exists approximately 55\% of the variance that does not have an
experimental cause.
The composite-retrieval score resulted in a higher correlation with
semantic similarity ($r{=}0.849$), confirming the assumption that
utilizing a multi-metric approach reduces the vulnerability of
relying on a single metric.

\subsection{Illustrative Examples}
\label{sec:cases}

\paragraph{Example 1: High P@K retrieval, Low RCA.}
SCTP Failure (S02): The Hybrid+Rerank model retrieved all files
relevant for locating data related to this case; however, the LLM
was not able to identify the whole chain of sequentially affected
data specifically referenced in the completion of this case.

\paragraph{Example 2: Low P@K retrieval, High RCA.}
Handover failure (S03): Although the Dense-Only model missing RAN
entries produced an invalid P@K score, it produced an accurate
chained reasoning for this case based on previous experiences and
knowledge within the domain.

\paragraph{Example 3: Iteration was found to be beneficial.}
On failure of Authentication (S07): Initially broad first-pass
retrieval; however, during query refinement, injecting the terms
``SQN'' and ``AUSF'' resulted in retrieving the proper evidence on
the 2nd iteration.

\paragraph{Example 4: Iteration was detrimental.}
For scheduler starvation (S08) Hybrid+Iter.: An incorrect initial
examination of the data produced drift during the query leading to
the retrieval of unrelated logs from a transport log.

% ═══════════════════════════════════════════════════════════════════════════════
% 6. DISCUSSION
% ═══════════════════════════════════════════════════════════════════════════════

\section{Discussion}
\label{sec:discussion}

\subsection{Reasons for Retrieval-Enhanced Results Not Delivering on Reasoning}

The reasoning/retrieval disconnect results from different potential
failure modes that limit retrieval (i.e., the limitations of index
coverage, fusion weights, and reranker accuracy) versus reasoning
(i.e., the limitation of an LLM to chain multiple sources of evidence
together and distinguish between relevant/not-relevant context).

The reasons for this disconnect can be summarized in three ways:
(i)~Precision@K is reasoning and retrieval agnostic.
Therefore, a document could be deemed to be relevant based on a
subject matter expert's annotation, but would not necessarily pass
through an LLM as part of its evidence chain when applying
multi-source reasoning.
(ii)~The task of evaluating end-task accuracy does not distinguish
between a correct answer that is based on LLM prior knowledge versus
correct based on retrieved evidence.
(iii)~When aggregating results across multiple queries, the overall
effectiveness of retrieval is reduced due to query drift.
For example, two different queries that exhibit opposite failed
retrieval states produce the same aggregated outcome; however, one
query will have a high return while another will have low returns.
Examples of this are illustrated in the scatter plots located in our
supplementary materials (Figure~1).
As indicated, of the total of 390 amounts of data collected at each
query configuration, there exists very little slope in the
corresponding regression ($r{=}0.244$), but rather a large degree of
vertical separation in the variance of what a user could be presented
with would indicate that regardless of what was returned as being a
relevant document to the query, there is significant variability in
the LLM's reasoning capabilities.

\subsection{RAG Evaluation Impact on Multi-Hop Retrieval}

To properly evaluate multi-hop RAG, these evaluations should:
(i)~separate retrieval and reasoning from one another;
(ii)~hold LLMs constant while varying retrieval;
(iii)~provide per-query performance metrics; and
(iv)~provide metrics of semantic accuracies based on more than simple
keyword overlap.
The mentioned are also applicable to multi-hop QA, RAG for scientific
integration, code assistance, and agentic retrieval.

\subsection{Multi-Mode RAG Impact on Retrieval/Reasoning Gap}

The retrieval vs.\ reasoning gap will likely be worsened when applied
to multi-modal environments.
When evidence of heterogeneous type (i.e., image, table, video, and
written evidence) from multiple sources is retrieved, there is no
guarantee that each of the multi-modal retrievals will be
successfully synthesized into coherent reasoning by a generation
model, even if retrieval has very high recall rates for each mode.
As noted above concerning how difficult it is to effectively
construct a reasoning framework from disparate modes of evidence, the
multi-hop research that we are studying provides relevant multi-hop
evidence types, but since each of these types are not always
modularized by construction, a retrieval model must also provide a
proper and aligned conceptual construct across each type of evidence.
The construct comparisons between varied retrieval types and a model
that performs the actual reasoning should provide a means to evaluate
multi-modal generation across large numbers of iterations.

\subsection{Adaptive to Fixed Iteration}

The adaptive stopping sequence provides an efficient means of
obtaining LLM evidence, while allowing for temporary downtimes in
confidence.
However, systems utilizing fixed iteration processes will have drift
or change in user queries (see Example~4) and the first-pass outcome
of a fixed system will affect the overall return due to the absence
of precision-oriented reranking.

% ═══════════════════════════════════════════════════════════════════════════════
% 7. LIMITATIONS AND FUTURE WORK
% ═══════════════════════════════════════════════════════════════════════════════

\section{Limitations and Future Work}
\label{sec:limitations}

\textbf{Text-only Evaluation.}
All experiments use only text as input for retrieval and generation.
More work is needed to investigate whether the retrieval/reasoning
gap still exists or increases when multimodal retrieval-augmented
generation (RAG) is used (for example, with image+text or video+text
evidence).

\textbf{Synthetic Dataset.}
The corpus used in this study was synthetically generated, with known
ground truth.
While this allows for controlled assessment, there are many factors
that may differ between the synthetic and true noise distributions,
as well as the level of uncertainty related to the annotation and the
size of the production systems.
Therefore, it will be necessary to conduct more work using real-world
telecommunications logs and established multi-hop benchmark datasets
(e.g., HotpotQA and MuSiQue) in order to assess the
generalizability of the models developed in this paper.

\textbf{LLM-as-Judge Bias.}
The development of cross-model judges and human evaluation will
provide additional support for the automated scoring methodologies.

\textbf{Single Domain.}
Performing the same experiment using other domains (e.g.,
code-assistant, scientific Q\&A, and medical) will strengthen the
generality claim from this work.

\textbf{Confidence Calibration.}
The composite confidence score is only a heuristic, not a formally
calibrated metric.

Future work will consist of: (i)~multimodal retrieval--reasoning
experiments that include image+text evidence; (ii)~employing
larger-than-individual-size corpora with bootstrap significance
testing; (iii)~performing a human-calibrated evaluation of these
models; (iv)~adapting domain-specific embeddings created using
contrastive learning methods; (v)~developing graph-based systems for
causal retrieval; and (vi)~fine-tuning the reader models to
determine whether such fine-tuning allows for advancement in
reasoning performance and reduction of the retrieval/reasoning gap.

% ═══════════════════════════════════════════════════════════════════════════════
% 8. CONCLUSION
% ═══════════════════════════════════════════════════════════════════════════════

\section{Conclusion}

The retrieval--reasoning gap has been demonstrated in this work:
across six different ablation configurations and the sixty-five
queries, P@K varied 0.24 ($p{<}0.001$) while RCA\textsubscript{kw}
only varied by 0.05 ($p{=}0.08$, n.s.) with Pearson $r{=}0.244$
between retrieval and reasoning.
In our investigation of multiple LLMs with fixed retrieval and their
corresponding reasoning accuracies, we show that, at times, simply
because of the model capability they displayed there was a reasoning
accuracy difference of as much as 0.103.
Of the forty-three percent of our per-query analysis, we found all
cases fall into the off-diagonal area of the retrieval--reasoning
space indicating that retrieval and reasoning do not agree upon the
same solution.
We also introduced an approach to break down (decompose) the LLMs to
create a single evaluation (composite metric) that showed good
agreement ($r{=}0.849$ with semantic evaluation) with the original
semantic evaluation data.
Therefore, the implication is clear: though both retrieval and
reasoning improvements must be pursued, improving retrieval alone is
not enough for multi-hop RAG success.
Thus, reasoning quality must be independently addressed and also
measured during the evaluation process of RAGs.

% ═══════════════════════════════════════════════════════════════════════════════
% REFERENCES
% ═══════════════════════════════════════════════════════════════════════════════

\bibliographystyle{plainnat}

\begin{thebibliography}{15}

\bibitem[Asai et~al.(2024)]{asai2024selfrag}
Asai, A., Wu, Z., Wang, Y., Sil, A., and Hajishirzi, H.
\newblock Self-RAG: Learning to retrieve, generate, and critique
through self-reflection.
\newblock In \emph{ICLR}, 2024.

\bibitem[Gao et~al.(2023)]{gao2023ragsurvey}
Gao, Y., Xiong, Y., Dibia, V., et~al.
\newblock Retrieval-augmented generation for large language models: A
survey.
\newblock \emph{arXiv preprint arXiv:2312.10997}, 2023.

\bibitem[Jiang et~al.(2023)]{jiang2023activerar}
Jiang, Z., Xu, F.~F., Gao, L., et~al.
\newblock Active retrieval augmented generation.
\newblock In \emph{EMNLP}, 2023.

\bibitem[Khattab and Zaharia(2020)]{khattab2020colbert}
Khattab, O. and Zaharia, M.
\newblock ColBERT: Efficient and effective passage search via
contextualized late interaction over BERT.
\newblock In \emph{SIGIR}, 2020.

\bibitem[Lewis et~al.(2020)]{lewis2020rag}
Lewis, P., Perez, E., Piktus, A., et~al.
\newblock Retrieval-augmented generation for knowledge-intensive NLP
tasks.
\newblock In \emph{NeurIPS}, 2020.

\bibitem[Nguyen et~al.(2016)]{nguyen2016msmarco}
Nguyen, T., Rosenberg, M., Song, X., et~al.
\newblock {MS MARCO}: A human generated machine reading comprehension
dataset.
\newblock In \emph{NeurIPS Workshop on Cognitive Computation}, 2016.

\bibitem[Nogueira and Cho(2019)]{nogueira2019passage}
Nogueira, R. and Cho, K.
\newblock Passage re-ranking with BERT.
\newblock \emph{arXiv preprint arXiv:1901.04085}, 2019.

\bibitem[Reimers and Gurevych(2019)]{reimers2019sentencebert}
Reimers, N. and Gurevych, I.
\newblock Sentence-BERT: Sentence embeddings using Siamese
BERT-networks.
\newblock In \emph{EMNLP}, 2019.

\bibitem[Robertson and Zaragoza(2009)]{robertson2009bm25}
Robertson, S. and Zaragoza, H.
\newblock The probabilistic relevance framework: {BM25} and beyond.
\newblock \emph{Foundations and Trends in IR}, 3(4):333--389, 2009.

\bibitem[Thakur et~al.(2021)]{thakur2021beir}
Thakur, N., Reimers, N., et~al.
\newblock {BEIR}: A heterogeneous benchmark for zero-shot evaluation
of information retrieval models.
\newblock In \emph{NeurIPS Datasets and Benchmarks}, 2021.

\bibitem[Wang et~al.(2023)]{wang2023queryrewrite}
Wang, L., Yang, N., and Wei, F.
\newblock Query2doc: Query expansion with large language models.
\newblock In \emph{EMNLP}, 2023.

\bibitem[Es et~al.(2024)]{es2024ragas}
Es, S., James, J., Espinosa-Anke, L., and Schockaert, S.
\newblock {RAGAS}: Automated evaluation of retrieval augmented
generation.
\newblock In \emph{EACL (System Demonstrations)}, 2024.

\bibitem[Saad-Falcon et~al.(2024)]{saad2024ares}
Saad-Falcon, J., Khattab, O., Potts, C., and Zaharia, M.
\newblock {ARES}: An automated evaluation framework for
retrieval-augmented generation systems.
\newblock In \emph{NAACL}, 2024.

\bibitem[Upadhyay et~al.(2024)]{upadhyay2024trecrag}
Upadhyay, S., Pradeep, R., Thakur, N., et~al.
\newblock A large-scale study of relevance assessments with large
language models: An initial look at the {TREC 2024 RAG Track}.
\newblock \emph{arXiv preprint arXiv:2411.08275}, 2024.

\bibitem[Liu et~al.(2024)]{liu2024lostmiddle}
Liu, N.~F., Lin, K., Hewitt, J., et~al.
\newblock Lost in the middle: How language models use long contexts.
\newblock \emph{TACL}, 12:157--173, 2024.

\end{thebibliography}

\end{document}
